\appendix

\section{Additional Analytical Details}\label{app:analysis}

\subsection{Quadratic coalition objectives and private travel time}

This subsection clarifies the economic normalization behind the quadratic specification. The integral in \eqref{eq:coalition-objective} supplies the common marginal road cost. A quadratic term is then needed if the coalition is meant to internalize the delay generated by its own flow rather than merely take that delay as given.

For affine costs, the physical part of coalition $C$'s objective on edge $e$ is
\[
 \int_0^{F_e^C}[a_e+b_e(x_e^{-C}+s)]\,ds
 =a_eF_e^C+b_ex_e^{-C}F_e^C+\frac12b_e(F_e^C)^2.
\]
Coalition travel time on that edge is
\[
 F_e^C[a_e+b_e(x_e^{-C}+F_e^C)]
 =a_eF_e^C+b_ex_e^{-C}F_e^C+b_e(F_e^C)^2.
\]
Thus the management penalty $\frac12b_e(F_e^C)^2$ converts the physical integral into the coalition's private total travel time on that edge. The equality is a normalization for the affine case, not a claim that every application must use total travel time as its private objective. For a singleton, $F_e^C=f_e^i$. For a fully integrated identity-preserving coalition, the governance matrix $b_e\mathbf 1_C\mathbf 1_C^\top$ produces exactly this aggregate square. A diagonal matrix retains separate member objectives; off-diagonal elements add cross-member internalization. The general quadratic form in \eqref{eq:quadratic-objective} allows these terms to coexist with route-specific priorities $q_C$.

\subsection{Projected inverse in a null-space basis}

Let $Z_C$ have orthonormal columns spanning $\ker\Gamma_C$. Any feasible active flow is $y_C^0+Z_Cu$. The coalition's fixed-support quadratic first-order condition restricted to this tangent is
\[
 (Z_C^\top H_CZ_C)u
 =-Z_C^\top[H_Cy_C^0+\Lambda_C^\top(a+Bx)+q_C].
\]
The tangent positive-definiteness condition \eqref{eq:coalition-tangent-pd} makes the coefficient invertible and gives
\[
 P_C=Z_C(Z_C^\top H_CZ_C)^{-1}Z_C^\top.
\]
This formula remains valid when $H_C$ is singular on route-space directions that cannot be generated by an OD-preserving movement. Exact duplicate routes must be combined or quotiented out when they create a zero-curvature vector in $\ker\Gamma_C$. The condition is therefore a statement about curvature on feasible deviations, not a requirement that the primitive objective be strictly convex in every route coordinate.

\subsection{HDV projection identities}

The HDV block enters the theorem through feasible reallocation directions, not through an additional management objective. The matrix $\mathcal A_H$ removes the edge-flow component that HDVs can generate while preserving their OD totals and equal-cost condition.

Let $K_H=L_H^\top BL_H$. Under \eqref{eq:hdv-tangent-pd},
\[
 \mathcal A_H=I-L_HK_H^{-1}L_H^\top B.
\]
The identities in Lemma~\ref{lem:screening} follow from
\[
 \mathcal A_HL_H=L_H-L_HK_H^{-1}K_H=0
\]
and
\[
 \mathcal A_H^2
 =I-2L_HK_H^{-1}L_H^\top B
 +L_HK_H^{-1}K_HK_H^{-1}L_H^\top B
 =\mathcal A_H.
\]
If $\mathcal A_Hv=0$, then $v=L_HK_H^{-1}L_H^\top Bv$, proving $\ker\mathcal A_H=\range(L_H)$. Symmetry of $B$ and $K_H$ gives $\mathcal A_H^\top B=B\mathcal A_H$.

\subsection{BPR continuation diagnostic}

The fixed-reference BPR implementation used in the main experiments is a convex potential-based behavioral continuation model; it is not the exact nonlinear own-travel-time atomic game unless the physical slope is constant.

The exact activated-response theorem is affine. It should not be read as a closed-form solution for the BPR continuation game. For BPR costs, the monotonicity proof in Proposition~\ref{prop:coalition-monotone} applies directly because BPR costs are continuous and nondecreasing and the coalition management matrices are PSD. The earlier perceived-slope specification has a more restrictive pointwise symmetric-Jacobian condition. On one edge, with fleet flows $f^k$, coefficients $\beta_k$, and $x=\sum_kf^k+h$, the condition is
\begin{equation}\label{eq:bpr-exact-psd}
 \frac{c''(x)}{c'(x)}
 \sqrt{\sum_k\beta_k(f^k)^2}\leq2.
\end{equation}
For a BPR exponent $\eta>1$, a readily checked sufficient bound is
\begin{equation}\label{eq:bpr-coarse-psd}
 \beta_{\max}\sigma\leq\frac{4}{(\eta-1)^2},
 \qquad \sigma=\max_k f^k/x.
\end{equation}
At $\eta=4$ the threshold is $4/9$. A pointwise pass records membership in a sufficient region; a fail removes that guarantee but is not evidence of multiplicity. All BPR profiles in the paper are certified independently by VI residual and omitted-route tests. Those numerical certificates establish approximate equilibrium of the solved route-generation problem; they do not extend the affine signature theorem to nonlinear BPR costs.

\subsection{Symbolic check of the canonical regime derivative}

For the bottleneck theorem, substituting
$x_1=(H+\Sigma x^{UE})/(1+\Sigma)$ into
$J=Na+t_1(x_1^2-x^{UE}x_1)$ and differentiating gives
\[
 \frac{dJ}{d\Sigma}
 =t_1(2x_1-x^{UE})\frac{x^{UE}-H}{(1+\Sigma)^2}.
\]
The identity
\[
 2x_1-x^{UE}=\frac{2H+(\Sigma-1)x^{UE}}{1+\Sigma}
\]
then yields \eqref{eq:bottleneck-tstt-derivative}. The repository script
\texttt{verify\_regime\_theorem.py} checks this simplification symbolically
and records the zero residual and threshold in
\texttt{regime\_theorem\_verification.json}. The script is a reproducibility
check; the proof in Theorem~\ref{thm:canonical-regimes} supplies the
mathematical argument.

\section{Sioux Falls Data and Numerical Certification}\label{app:sf-params}

Sioux Falls has $24$ nodes, $76$ directed links, and $528$ positive-demand OD pairs. Demand, free-flow times, and capacities are read from the standard network files in the repository. Total demand is $360{,}600$. Physical costs use \eqref{eq:bpr}. The all-HDV user-equilibrium and system-optimum TSTTs are $7{,}480{,}226$ and $7{,}194{,}256$, respectively.

The initial route set contains the three shortest free-flow routes per OD. It is expanded by player-specific shortest-path column generation, because a route that is unattractive at free flow can become attractive under the marginal objective induced by coalition governance. A route is added when its current marginal route cost is below the minimum held-route marginal cost by more than $10^{-6}$ times the route-cost scale. The VI is solved with adaptive Korpelevich extragradient and Euclidean projection onto each identity-preserving company--OD simplex. The accepted profiles satisfy normalized VI gap at most $10^{-6}$ and full-network omitted-route slack at most $10^{-6}$.

The composition experiment contains $40$ profiles for four equal-mass
companies with objective types $(0.5,0.5,1.5,1.5)$. It crosses two penetration
values, five coordination values, and four partitions. The degree experiment
contains $54$ profiles across $n\in\{2,4,8\}$, two penetration values, three
coordination values, and all declared balanced partitions.

The sensitivity experiment contains $180$ profiles crossing three objective gaps, three reference-slope scales, five coordination values, two penetrations, and the assortative/mixed partition pair. The complete $n=4$ partition scan contains all $15$ set partitions at two positive coordination values and two penetrations, for $60$ profiles. The pooling decomposition contains four profiles comparing identity-preserving and operational-pooling feasibility. All $338$ Sioux Falls coalition-extension profiles are certified.

The cross-network audit contains a five-edge Braess topology and a directed
$3\times3$ grid with three positive-demand ODs. For each topology, the design
crosses four AV penetration levels, four governance intensities, and all $15$
partitions of four companies, for $480$ profiles. All profiles are certified;
the maximum normalized VI gap is $9.971\times10^{-7}$ and the maximum relative
omitted-route slack is $3.79\times10^{-16}$. The raw profile file is
\texttt{cross\_network\_regime\_results.json}; the state and policy summaries
are in \texttt{cross\_network\_regime\_summary.json}.

Policy regret is computed after, rather than during, equilibrium solution.
\nolinkurl{coalition_design_analysis.py} calculates set-valued Sioux Falls
optima, one- and two-step merge paths, fixed-policy regret, and the
governance-cost lower envelope. \nolinkurl{cross_network_policy_analysis.py}
applies the same local policies to the two independent topologies. The
validator \nolinkurl{validate_regime_package.py} checks profile counts,
certification thresholds, finiteness, Sioux Falls policy consistency, and the
reported cross-network regret values.

The fixed-HHI spatial experiment contains $82$ certified profiles. Four equal-mass companies have HHI $0.25$ at both $\alpha=0.5$ and $0.9$; iterative proportional fitting preserves every aggregate OD row and every fleet-mass column while varying spatial OD portfolios. The response-decomposition experiment contains $600$ certified rows, the penetration scan $500$, and the targeted BPR audit $10$. The raw and summary JSON files, scripts, and validators are included in the repository.

\begin{table}[htbp]
\centering
\scriptsize
\setlength{\tabcolsep}{3pt}
\caption{Numerical certification summary. All rows satisfy both acceptance criteria. Exact PSD entries are pointwise diagnostics for the earlier perceived-slope family, not acceptance criteria for the new coalition quadratic family.}
\label{tab:certificates}
\begin{tabular}{lcccr}
\toprule
Experiment & certified & max VI gap & max relative oracle slack & rows\\
\midrule
coalition composition & $40/40$ & $1.0\times10^{-6}$ & $1.0\times10^{-6}$ & 40\\
coalition degree & $54/54$ & $9.97\times10^{-7}$ & $9.91\times10^{-7}$ & 54\\
coalition sensitivity & $180/180$ & $1.0\times10^{-6}$ & $1.0\times10^{-6}$ & 180\\
all $n=4$ partitions & $60/60$ & $1.0\times10^{-6}$ & $1.0\times10^{-6}$ & 60\\
Braess/grid regime audit & $480/480$ & $9.97\times10^{-7}$ & $3.79\times10^{-16}$ & 480\\
identity/pooling decomposition & $4/4$ & $9.84\times10^{-7}$ & $3.2\times10^{-8}$ & 4\\
fixed-HHI spatial OD & $82/82$ & $1.0\times10^{-6}$ & $1.0\times10^{-6}$ & 82\\
response decomposition & $600/600$ & $1.0\times10^{-6}$ & $9.0\times10^{-7}$ & 600\\
penetration scan & $500/500$ & $1.0\times10^{-6}$ & $4.2\times10^{-7}$ & 500\\
BPR anchor audit & $10/10$ & $1.0\times10^{-6}$ & $1.2\times10^{-7}$ & 10\\
\bottomrule
\end{tabular}
\end{table}

\section{Reproducibility and Scope Ledger}\label{app:reproducibility}

The current coalition solver extends the previous fleet solver through optional coalition governance arguments while preserving the earlier model as a default nested case. The coalition objective uses fixed reference curvature $\bar b_e=c'_e(x_e^{UE})$ and the PSD matrix in \eqref{eq:gamma-governance}. The company-count scan uses equal company masses within each $n$ and holds total AV mass fixed across $n$. Singleton profiles at nonzero $\gamma$ are generated from the analytically verified partition invariance at $\gamma=0$; they are labeled as derived rows in the raw JSON and are not counted as independent solver runs. These implementation choices define the declared counterfactual menu; they do not estimate how firms would choose a coalition in an unfixed market.

The empirical claims in Section~\ref{sec:experiments} are limited to the declared Sioux Falls benchmark and governance design. The objective types are controlled primitives, not estimates from operator data. Operational pooling, endogenous coalition stability, transfers, pricing, passenger choice, and platform welfare remain outside the reported core results. The theoretical claims are likewise conditional: existence and monotonicity use continuous nondecreasing physical costs, whereas the activated signature and its minimality use affine costs, fixed supports, and the stated tangent-rank conditions.

The repository includes the extension and cross-network scripts, raw and
summary JSON outputs, symbolic checks, fail-closed validators, and
figure-generation code. Pooling is implemented as a coalition-level simplex
over member-route coordinates; its comparison with identity-preserving
coordination is interpreted only under the proportional OD and common-route
design reported in Section~\ref{sec:experiments}. The local merge policies use
full scans only as audit baselines for regret. Their zero observed two-step
regret is not encoded as an assumption and is not claimed as a general
optimality theorem.
